\section{Implementation and validation}
\label{sec:impl}

\subsection{Software design}

The strategy is implemented as a Python package (\code{pca\_stat\_arb}) built on NumPy and
pandas. \Cref{tab:modules} lists its modules. Three design decisions matter for the
correctness and extensibility of the results.

\begin{enumerate}
  \item \textbf{A single data contract.} Every data source implements one interface,
        \code{ReturnsSource}, whose output is a table of daily total returns indexed by date and
        security. Nothing downstream depends on where the data came from, so replacing Yahoo
        Finance with CRSP requires only a new loader.
  \item \textbf{Configuration as data.} All parameters live in typed, immutable configuration
        objects loaded from YAML files and are checked for validity when loaded. Each experiment
        is a named set of overrides of the baseline configuration, so every reported variant can
        be reproduced from a file.
  \item \textbf{Vectorised daily step.} Within each day all computations run on whole arrays
        rather than looping over stocks: one symmetric eigendecomposition, one multi-target
        least-squares solve for all regressions in \cref{eq:regression}, and closed-form AR(1)
        estimates for all stocks at once. A full backtest over 15 years and 500 stocks takes
        about 45 seconds on a laptop, which makes systematic parameter sweeps practical.
\end{enumerate}

\begin{table}[ht]
\centering
\caption{Package structure.}
\label{tab:modules}
\small
\begin{tabular}{@{}l>{\raggedright\arraybackslash}p{0.52\textwidth}@{}}
\toprule
Module & Responsibility \\
\midrule
\code{config} & typed configuration, validation \\
\code{data.base}, \code{data.yfinance\_source} & data interface; Yahoo Finance loader with caching and cleaning \\
\code{data.universe} & constituent list snapshot \\
\code{pca} & eigenportfolios (\cref{sec:method:pca}) \\
\code{residuals} & factor regression, OU fit, s-scores (\cref{sec:method:regression,sec:method:ou,sec:method:sscore}) \\
\code{signals} & trading rule (\cref{sec:method:rule}) \\
\code{backtest} & daily loop, hedging, costs (\cref{sec:method:portfolio,sec:method:costs,sec:method:accounting}) \\
\code{metrics} & performance statistics, in-sample/test split \\
\code{scripts/experiments.py} & parallel variant runner, in-sample by default \\
\code{scripts/make\_report\_assets.py} & all figures, tables and numbers in this report \\
\bottomrule
\end{tabular}
\end{table}

\paragraph{Reproducibility of this report.}
Every figure, results table and quoted statistic in this report is produced by
\code{scripts/make\_report\_assets.py} directly from the code's output. The numbers in the text are
inserted as generated \LaTeX{} macros rather than typed by hand, so the report cannot drift out
of sync with the code. The script truncates the data at the split date before doing anything
else, so no held-out data can enter the report until the final evaluation.

\subsection{Validation on synthetic data}
\label{sec:impl:tests}

Backtests can fail silently: a subtle bug can produce plausible-looking but wrong P\&L. The
implementation is therefore tested on a simulated market where the true structure is known.
Returns are generated from three factors (a dominant market factor and two zero-mean
``sector'' factors) plus idiosyncratic residuals whose cumulative sums follow an OU process with
known $\kappa$. \Cref{tab:tests} summarises the automated tests, all of which pass.

\begin{table}[ht]
\centering
\caption{Automated tests. All run on synthetic data with known structure.}
\label{tab:tests}
\begin{tabular}{p{0.38\textwidth}p{0.54\textwidth}}
\toprule
Test & Property checked \\
\midrule
Spectrum vs.\ Marchenko--Pastur edge & exactly the 3 true factors have eigenvalues above $(1+\sqrt{N/T})^2$; the first eigenvector is the market \\
Orthogonality & regression residuals are orthogonal to the factors, and $X_{iL}=0$ \\
OU recovery & the AR(1) mapping recovers the true $\kappa$ (up to known small-sample bias) \\
s-score eligibility & slow or non-mean-reverting residuals yield no s-score \\
Trading rule & every transition in \cref{tab:rule}, including same-day reversal \\
\textbf{No look-ahead} (three variants) & replacing all data after day $t$ with random noise leaves every holding up to day $t$ unchanged, for daily hedges, frozen hedges with volume adjustment, and the no-trade band \\
End-to-end profitability & the strategy is profitable on the synthetic mean-reverting market \\
Hedge modes & frozen hedges give identical signals with lower turnover, and are released on exit \\
No-trade band & every executed trade exceeds $\delta$; turnover falls \\
Costs & costs reduce net P\&L and leave gross P\&L unchanged \\
Missing data & a stock with a data gap is not traded until it has a full window again \\
\bottomrule
\end{tabular}
\end{table}

The look-ahead tests are the most important. Look-ahead bias, the use of information that would
not have been available at the time of the trade, is the most common way a backtest overstates
performance. Rather than relying on a careful reading of the code, these tests check the
property mechanically for every hedging mode.

\paragraph{Consistency checks on real data.}
On the real data we additionally verified that the recorded states never violate the trading
rule. Across all stocks and in-sample days, no flat stock with $|s_i| > \bar s_{\text{open}}$
failed to open the corresponding position.
